
où $p$ est un
morphisme lisse. On peut de plus supposer $\underline{\Omega}^1_{V/Y}$, $\underline{\Omega}^1_{V'/Y}$, $\underline{\Omega}^1_{V'/V}$ de rang constant, ainsi que les faisceaux conormaux de $i$ et $j$.

Du diagramme commutatif (où tous les morphismes sont lisses) :
\[
\begin{tikzcd}[column sep=large,row sep=large]
V' \arrow[rr,"p"] \arrow[dr] && V \arrow[dl]\\
& Y &
\end{tikzcd}
\]
on déduit :
\begin{equation}\tag{1.8.1}
\operatorname{rang}(\underline{\Omega}^1_{V'/Y})=
\operatorname{rang}(\underline{\Omega}^1_{V/Y})+
\operatorname{rang}(\underline{\Omega}^1_{V'/V}).
\end{equation}

Soit $X'=V'\times_V X$; alors le rang de $\underline{\Omega}^1_{V'/V}$ est égal à celui de $\underline{\Omega}^1_{X'/X}$. Par ailleurs, $j$ définit une section $X\xrightarrow{j'} X'$, qui est une immersion régulière de codimension égale au rang de $\underline{\Omega}^1_{X'/X}$ (EGA IV 17.10.4). Si $i'$ est l'immersion déduite de $i$ par le changement de base $V'\to V$, on a $j=i'\circ j'$ et par suite (VII 1.7) :
\begin{equation}\tag{1.8.2}
\operatorname{rang}(\underline{N}_{V'/X})=
\operatorname{rang}(\underline{N}_{V/X})+
\operatorname{rang}(\underline{\Omega}^1_{V'/V}).
\end{equation}

En comparant (1.8.1) et (1.8.2), on en déduit la proposition.

\textbf{Définition 1.9.} \uline{L'entier} $d_x$ \uline{de la proposition} 1.8 \uline{s'appelle} dimension relative virtuelle \uline{de} $X$ \uline{sur} $Y$ (\uline{ou de} $f$) \uline{au point} $x$.

La dimension relative virtuelle de $X$ sur $Y$ est donc une fonction localement constante sur $X$, à valeurs entières (positives ou négatives). Si $f$ est un morphisme lisse de dimension relative $n$, sa dimension relative virtuelle est $n$. Si $f$ est une immersion régulière de codimension $n$, sa dimension relative virtuelle est $-n$.

\textbf{Proposition 1.10.} \uline{Soient} $f:X\to Y$ \uline{et} $g:Y\to Z$ \uline{deux morphismes d'intersection complète}, \uline{et soit} $x$ \uline{un point de} $X$. \uline{Alors} :
\[
\operatorname{dim.rel.virt.}_x(g\circ f)=
\operatorname{dim.rel.virt.}_x(f)+
\operatorname{dim.rel.virt.}_{f(x)}(g)
\]
